\documentclass[border=4pt]{standalone}
\usepackage{amsmath,amssymb,mathtools,xcolor,varwidth}
\definecolor{circ}{HTML}{2ea043}
\definecolor{curve}{HTML}{1f6feb}
\definecolor{excess}{HTML}{cf3b2f}
\begin{document}
\begin{varwidth}{28em}
\large\bfseries
Finally Newton circumscribes \textcolor{circ}{the parallelograms $aKbl$, $bLcm$, $cMdn$},
each rising above \textcolor{curve}{the curve} to its upper vertex.
The difference between \textcolor{circ}{the circumscribed figure} and the inscribed figure
is the sum of the little tops $Kl$, $Lm$, $Mn$, $Do$; since all bases are equal,
this sum is \textcolor{excess}{the single rectangle $ABla$}, of breadth $AB$ and height $Aa$.
But as $AB$ is diminished $in\ infinitum$, \textcolor{excess}{this rectangle} becomes less
than any given space, so by Lemma~I the inscribed and \textcolor{circ}{circumscribed} figures
become ultimately equal, and $much\ more$ is \textcolor{curve}{the intermediate curvilinear figure}
equal to either.
$Q.E.D.$
\end{varwidth}
\end{document}
